\section{The July run}\label{app:july}
\setcounter{table}{0}\setcounter{figure}{0}\setcounter{equation}{0}

The July run (\cref{tab:runs}), pre-registered before it was executed, with the weaker
public record of \cref{sec:methods:prereg} (\ref{app:prereg}), trained the same kinds of
network as \cref{sec:results2d} on the same corpus, split and seeds, with the earlier code,
which did not multiply the output by $F_{\max}$ (\cref{sec:methods:networks}). Its main
results were first reported in the companion note \cite{CaoSong2026}. It also trained what
the two-dimensional comparison did not:
the supervised transformer with the energy term of \cref{sec:methods:training}, the
supervised transformer trained on $\mathcal{L}_D$ from the label-free transformer's
parameters, and the graph network with every number of labelled instances. It trained no
stiffness-norm transformer.

\subsection{Accuracy and label efficiency}\label{app:july:accuracy}

% generated by scripts/make_cmame_material.py from the records; do not edit
\begin{table}[tbp]
\centering
\footnotesize
\setlength{\tabcolsep}{4pt}
\caption{Two-dimensional plates, the July run with the earlier code: every metric at the largest label budget, on the instances of \cref{tab:2d}.}
\label{tab:2djuly}
\resizebox{\ifdim\width>\textwidth\textwidth\else\width\fi}{!}{%
\begin{tabular}{lcccccc}
\toprule
Model & \shortstack{Displacement \\ error} & \shortstack{Relative \\ energy gap} & \shortstack{von Mises \\ error} & \shortstack{Peak stress \\ error} & \shortstack{Critical-region \\ recall} & \shortstack{Worse than \\ zero field} \\
\midrule
Label-free & 0.166 $\pm$ 0.002 & 0.0817 $\pm$ 0.0040 & 0.210 $\pm$ 0.003 & 0.228 $\pm$ 0.001 & 0.739 $\pm$ 0.003 & 3/768 \\
Supervised & 0.160 $\pm$ 0.003 & 0.418 $\pm$ 0.037 & 0.433 $\pm$ 0.031 & 0.334 $\pm$ 0.036 & 0.584 $\pm$ 0.027 & 52/768 \\
Supervised + energy term & 0.161 $\pm$ 0.005 & 0.166 $\pm$ 0.040 & 0.257 $\pm$ 0.007 & 0.246 $\pm$ 0.022 & 0.708 $\pm$ 0.012 & 14/768 \\
Label-free, then supervised & 0.166 $\pm$ 0.006 & 0.236 $\pm$ 0.032 & 0.336 $\pm$ 0.020 & 0.301 $\pm$ 0.019 & 0.645 $\pm$ 0.008 & 21/768 \\
Graph network, supervised & 0.167 $\pm$ 0.000 & 2.57 $\pm$ 0.25 & 1.10 $\pm$ 0.07 & 1.32 $\pm$ 0.13 & 0.322 $\pm$ 0.027 & 481/768 \\
\midrule
Nearest-neighbour field & 0.404 & 0.959 & 0.735 & 0.643 & 0.515 & 94/256 \\
Scaled polynomial & 1.42 & 873 & 11.4 & 30.8 & 0.159 & 256/256 \\
\midrule
Zero field & 1.00 & 1.00 & 1.00 & 1.00 & -- & -- \\
\bottomrule
\end{tabular}}

\smallskip
\begin{minipage}{\linewidth}\footnotesize\raggedright
256 validation instances, never trained on; four load cases each. Errors are relative to the reference finite-element solution, per load case, then averaged over the load cases; the relative energy gap of the zero field is 1. Trained networks: mean $\pm$ sample standard deviation over 3 seeds; the nearest-neighbour field and the scaled polynomial are deterministic.\par
Worse than zero field: instance-seed pairs whose relative energy gap, the mean over the four load cases as the evaluation records it, exceeds 1, that is, predictions further from the solution in the stiffness norm than the zero field on average; the count uses the reference solutions. The test of \cref{cor:zero} needs none and applies to single load cases (\cref{sec:methods:metrics}).\par
Critical-region recall: the share of the 10\% most stressed elements of the reference solution that are also among the 10\% most stressed of the prediction (higher is better; in every other column lower is better).\par
The July run (16 July 2026). Its code did not multiply the load scale back onto the output (\cref{sec:methods:networks}); we attribute the narrow spread of the displacement errors to this; no experiment isolated it.
\end{minipage}
\end{table}


\Cref{tab:2djuly} gives every metric with 1,024 labelled instances. The supervised
transformer had the lowest displacement error, \numTwoLabDispFour{} against
\numTwoFreeDispFour{} for the label-free transformer, which was \numTwoFreeDispExcessAbs{}
higher, and higher in each of the three seeds. On every other metric the label-free
transformer was the most accurate network in the table: its relative energy gap,
\numTwoFreeGap{}, was \numTwoLabOverFreeGap{} times lower than the supervised transformer's,
and its von Mises error (\numTwoFreeVm{} against \numTwoLabVm{}), peak stress error
(\numTwoFreePeak{} against \numTwoLabPeak{}) and critical-region recall (\numTwoFreeRecall{}
against \numTwoLabRecall{}) were better. Compared instance by instance within each seed, the
label-free transformer had the lower relative energy gap on \numPairTwoLabGap{} of the
\numTwoPairs{} instance-seed pairs and the lower von Mises error on \numPairTwoLabVm{}, but
the lower displacement error on only \numPairTwoLabDisp{}. The energy term lowered the
supervised transformer's relative energy gap to \numTwoAncGap{} and its von Mises error to
\numTwoAncVm{}, both still above the label-free values. The graph network matched the others
in displacement (\numTwoMgnDisp{}) but had a relative energy gap of \numTwoMgnGap{} and a
von
Mises error of \numTwoMgnVm{}: \numTwoMgnWorse{} of its \numTwoPairs{} instance-seed
predictions had a mean relative energy gap above one. The supervised transformer had
\numTwoLabWorse{} such predictions, the supervised transformer with the energy term
\numTwoAncWorse{}, and the label-free transformer \numTwoFreeWorse{}, one per seed.

The displacement errors of the five trained networks lay between \numTwoLabDisp{} and
\numTwoMgnDisp{}, a narrow band, while their relative energy gaps spanned a factor of
\numTwoGapSpread{}. We attribute the narrow band to the code: no network could know the
absolute amplitude of an instance's loads, which adds an error common to every displacement
and compresses their differences. With the later code, the displacement errors of the
label-free transformer, the supervised transformer and the graph network on the same
instances were \numJulyOverCmDispMin{} to \numJulyOverCmDispMax{} times lower
(\cref{tab:2d}), and those of the trained networks spanned a factor of \numCmDispSpread{}
rather than \numTwoDispSpread{}. This agrees with the attribution, but the two runs also
differ in their software versions, and no experiment isolated the output scaling.

% generated by scripts/make_cmame_material.py from the records; do not edit
\begin{table}[tbp]
\centering
\footnotesize
\setlength{\tabcolsep}{4pt}
\caption{The July run (two dimensions, earlier code): relative displacement error / relative energy gap against the number of labelled training instances (seed means).}
\label{tab:labeleffjuly}
\resizebox{\ifdim\width>\textwidth\textwidth\else\width\fi}{!}{%
\begin{tabular}{lcccc}
\toprule
Model & 16 & 64 & 256 & 1,024 \\
\midrule
Supervised & 0.272 / 6.26 & 0.193 / 2.72 & 0.176 / 1.26 & 0.160 / 0.418 \\
Supervised + energy term & 0.410 / 2.56 & 0.184 / 0.579 & 0.182 / 0.154 & 0.161 / 0.166 \\
Label-free, then supervised & 0.207 / 1.32 & 0.189 / 0.201 & 0.181 / 0.302 & 0.166 / 0.236 \\
Graph network, supervised & 0.684 / 92.6 & 0.425 / 80.6 & 0.181 / 7.55 & 0.167 / 2.57 \\
Nearest-neighbour field & 0.687 / 6.50 & 0.466 / 1.65 & 0.408 / 1.03 & 0.404 / 0.959 \\
Label-free (no labels) & 0.166 / 0.0817 & 0.166 / 0.0817 & 0.166 / 0.0817 & 0.166 / 0.0817 \\
\bottomrule
\end{tabular}}

\smallskip
\begin{minipage}{\linewidth}\footnotesize\raggedright
Seed means over 3 seeds; the nearest-neighbour field is deterministic, and the same as in \cref{tab:labeleff}. The graph network was trained at every budget in this run. The label-free row repeats its single value, which uses the same 1,024 instances without labels.
\end{minipage}
\end{table}


With no labels, the label-free transformer's relative energy gap was lower than every
supervised network's at every budget (\cref{tab:labeleffjuly}); its advantage over the
supervised transformer fell from \numTwoAdvMax{} with 16 labels to \numTwoAdvMin{} with
1,024. In displacement, the supervised transformer needed 1,024 labels to pass it. With 16
labels, the supervised transformer with the energy term was less accurate in displacement
than without it, and its per-seed errors scattered widely (\numTwoAncSixteenSeeds{}); the
diagnostic run later traced this to the instability of the gradient-scaled energy term with
very few labels (\ref{app:prereg}). With 16 labels the supervised transformer of the later
code was less accurate than that of the July run, with a relative energy gap of
\numCmLabGapSixteen{} against \numTwoLabGapSixteen{}.

\subsection{Pre-registered decisions}\label{app:july:decisions}

The gate of this run concerned the energy term and pretraining rather than the label-free
network alone, and all three of its conditions were met (\cref{tab:criteria}). Its first
condition read a separate energy-term sub-experiment of the same run, which trained its own
supervised transformers with and without the energy term in the same configuration: there,
the supervised transformer with the energy term was \numGOneSanityMin{} to
\numGOneSanityMax{} times more accurate in displacement than the zero field and beat the two
naive baselines and the plain polynomial (\cref{sec:methods:networks}), all fitted on 1,024
labelled instances, at every budget. The other two conditions read the networks of
\cref{tab:2djuly,tab:labeleffjuly}: with 64 labels, the energy term lowered the relative
energy gap by \numGOneGapReduction{} and the von Mises error by \numGOneVmReduction{}, and
label-free training followed by supervised training beat supervised training alone by
\numGOneFtGap{} in relative energy gap. The kill conditions on the label-free network were
not triggered: its displacement error was within the 30\% allowed by K2, and its relative
energy-gap advantage stayed above the 40\% floor at every budget. Two kill conditions were
triggered, and both claims were withdrawn: K5, on warm starts, which \cref{sec:cost:warm}
reports, and K6.

\paragraph{Cross-resolution regulariser (K6)}
A term that penalised differences between the latent representations of the same instance
meshed at two resolutions did not shrink the transfer gap, the increase of the
displacement error from the training resolution to a finer one; it widened it, by
\numCoarsenHighWiden{} for meshes coarsened by a factor of \numCoarsenHigh{} and by
\numCoarsenLowWiden{} for \numCoarsenLow{}. The regulariser plays no further part in this
paper.

\subsection{Further observations}\label{app:july:other}

\paragraph{Replication}
The diagnostic run (30 July 2026), which was not pre-registered, trained the label-free
transformer \numDiagSeeds{} more times with the same configuration, with seeds 0 to 5. The
first three repeat the seed numbers of the July run but did not reproduce its values
exactly; the diagnostic run ran on a different GPU model and a later release of the code
whose training path is unchanged. Over all \numTwoRuns{} trainings the displacement error
was \numTwoRunsDisp{} and the relative energy gap \numTwoRunsGap{} (mean and sample standard
deviation).

\paragraph{Cross-resolution transfer}
In the test of K6, label-free transformers trained for 100 epochs, with seed 0, on 512
coarsened meshes and without the regulariser were evaluated on the coarsened and on the
original meshes of 128 held-out instances. The displacement error rose from
\numCoarsenLowCoarse{} to \numCoarsenLowFine{} for meshes \numCoarsenLow{} times finer
than in training and from \numCoarsenHighCoarse{} to \numCoarsenHighFine{} for meshes
\numCoarsenHigh{} times finer. In this code version the output did not depend on the mesh
through $F_{\max}$; the contrast with three dimensions is taken up in
\cref{sec:transfer:amplitude}.
